\section{Capacity and one-sided smoothing}
\label{sec:rpi:orientation}

The mass decrease in conformal flow is controlled by the difference
between ADM mass and capacity. To show that this difference is
nonnegative, we first move a possibly singular frontier toward the
asymptotic end and smooth it with a definite mean-curvature sign.
The capacity does not decrease under this operation. We can then apply a
smooth mass--capacity inequality, proved below from positive mass.

Throughout this section $n\ge3$, $\omega=|S^{n-1}|$, and $\nu$
denotes the normal from a compact removed region toward the
asymptotically flat end. For a frontier $\Sigma$ we use
$H_\nu=\divg_\Sigma\nu$.
For a smooth compact conductor with exterior $(E,h)$, normalize
capacity by
\begin{equation}\label{eq:rpi:capacity}
 \mathfrak c(\Sigma,h)=\frac1{(n-2)\omega}
             \int_E|dw|_h^2\,dV_h,
 \qquad \Delta_h w=0,\quad w|_\Sigma=1,\quad w\to0.
\end{equation}
Equivalently, this is the normalized infimum of Dirichlet energy
among smooth compactly supported trial functions equal to one in a
neighborhood of the conductor. The same variational definition
applies to a nonsmooth compact conductor.
If one conductor contains another, its admissible class is smaller,
so its capacity is no smaller.

\subsection{A harmonic perturbation and its detached enclosure}

We first construct a boundary with a strict curvature sign; the next
subsection will smooth it. Let $(E,g)$ be a flow-slice exterior or
a smooth-ambient initial exterior as in Theorem~\ref{thm:rpi},
and assume its frontier is
nonempty. Let $w$ be its capacitary potential and put $v=w-1$.
The strong maximum principle gives
$-1<v<0$ in the open exterior, while $v=0$ on the frontier.
For $0<\varepsilon<1$ set
\[
 u_\varepsilon=\frac{1+\varepsilon v}{1-\varepsilon},\qquad
 \widetilde g_\varepsilon=u_\varepsilon^{4/(n-2)}g.
\]
Harmonicity gives
$R_{\widetilde g_\varepsilon}=u_\varepsilon^{-4/(n-2)}R_g\ge0$
on the exterior.
The perturbing factor tends to one at infinity, preserving the end
normalization of $g$. To apply the detachment
corollary, whose harmonic factor has boundary value one, express
the same metric as
\[
 g^{\rm base}_\varepsilon=(1-\varepsilon)^{-4/(n-2)}g,
 \qquad \widehat u_\varepsilon=1+\varepsilon v,
 \qquad
 \widetilde g_\varepsilon
   =\widehat u_\varepsilon^{4/(n-2)}g^{\rm base}_\varepsilon.
\]
Then $\widehat u_\varepsilon=1$ on the frontier and
$1-\varepsilon<\widehat u_\varepsilon<1$ in $E$.
If $g^{\rm base}_\varepsilon=\kappa^2g$, then
\[
 \Delta_{g^{\rm base}_\varepsilon}=\kappa^{-2}\Delta_g,
 \qquad \nu_{g^{\rm base}_\varepsilon}=\kappa^{-1}\nu_g,
 \qquad H^{g^{\rm base}_\varepsilon}=\kappa^{-1}H^g.
\]
Constant scaling preserves harmonicity, minimality and the local
almost-minimizing hypotheses, with changed constants. It also preserves
the signs and relative bounds of normal derivatives. At regular
boundary points the Hopf derivative of $\widehat u_\varepsilon$
in the direction $\nu$ is strictly negative, so the conformal
mean-curvature formula makes the original frontier a strict inner
barrier in $\widetilde g_\varepsilon$.

Apply the prescribed-curvature construction of
Section~\ref{sec:rpi:perturbed-candidate} to this fixed perturbation,
using its smooth forcing $h_0$, zero on the original filled side and
positive in its exterior. With the sufficiently small fixed volume
parameter $\delta>0$, Corollary~\ref{cor:rpi:flow-detachment} separates
the candidate from the full original frontier. Denote its connected
end exterior by $E'$. In the perturbed metric
$\widetilde g_\varepsilon$, a variation of its filled side by
$\varphi\nu$ has first variation
\begin{equation}\label{eq:rpi:first-variation}
 D\mathcal F[\varphi\nu]
       =\int_{\partial E'}(H_\nu+\delta h_0)\varphi\,dA,
\end{equation}
and therefore its regular frontier satisfies
\begin{equation}\label{eq:rpi:corrected-sign}
 H_\nu=-\delta h_0\le-\gamma<0.
\end{equation}
Here $\gamma>0$ exists because the entire detached frontier is a
compact subset of the original exterior, where $h_0>0$.
The perturbed metric is smooth on its neighborhood. All perturbation and
forcing parameters remain fixed through separation and smoothing;
no curvature or separation bound uniform in $\varepsilon$ is needed.

Strictness matters: if $u\equiv1$, the forcing is zero, and the
original frontier is outer minimizing and minimal, the original
exterior remains a minimizing competitor. Separation cannot follow
from $u\le1$ alone.

\begin{remark}[Orientation and the literature]
On the Euclidean exterior $r\ge a$, our convention gives
\[
 H_\nu=(n-1)/a>0,\qquad m_{\rm ADM}=0,\qquad
 w=(a/r)^{n-2},\qquad\mathfrak c=a^{n-2}>0.
\]
The smooth mass--capacity inequality therefore requires the
nonpositive sign with this normal. This example distinguishes the
statement here from the positive-$H_\nu$ inequality in
\cite[Lemma~3.11]{BiZhu2026}, which is version~1. Version~2 uses
the outward normal of the exterior in its Lemma~3.10, hence the
opposite normal at the inner frontier. Its Corollary~2.37 includes
the strict harmonic-factor hypothesis, and its Lemma~3.7 includes
the base-metric normalization \cite{BiZhu2026v2}.
\end{remark}

\subsection{Smooth frontiers with the required sign}
\label{sec:rpi:smoothing}

The distance-offset construction in
\cite[Section~3 and Proposition~2.6]{BiZhuSmoothing2026},
arXiv:2605.12403v2, supplies the geometric strategy for the next
proposition. The conclusion needed for capacity is conductor inclusion and a
mean-curvature bound. We prove precisely this conclusion using
Rifford's minimum-family theorem locally and a final smoothing of
transverse graphs. Boundary-area convergence is not needed.

\begin{proposition}[One-sided smoothing in a shrinking collar]
\label{prop:rpi:one-sided-smoothing}
Let $\Omega$ be an open subset of a smooth Riemannian manifold with
compact full boundary $\Sigma$. Assume that the metric is smooth on
an ambient neighborhood of $\Sigma$. Let $\Lambda\subset\Sigma$
consist of the points touched by a geodesic ball contained in $\Omega$.
Suppose that $\Lambda$ is a smooth embedded hypersurface, locally
bounds $\Omega$ on its accessible side, and has mean curvature at
least $\gamma>0$ for the normal pointing out of $\Omega$.

For every $\varepsilon>0$, there is an open set
$\Omega_\varepsilon\subset\Omega$ with smooth compact boundary such
that
\[
 \Omega\setminus\Omega_\varepsilon
 \subset\{x\in\Omega:\dist(x,\Sigma)<\varepsilon\},
 \qquad
 H_{\mathrm{out},\Omega_\varepsilon}\ge\gamma-\varepsilon.
\]
\end{proposition}

\begin{proof}
If $\Sigma$ is empty, take $\Omega_\varepsilon=\Omega$.
Otherwise we first choose a regular distance level, then take a short offset
of its superlevel, and finally smooth the resulting $C^{1,1}$
hypersurface. The offset controls curvature; the final smoothing
loses an arbitrarily small amount of that bound.

Write $\rho(x)=\dist(x,\Sigma)$ on $\Omega$. All distances and
geodesics below are taken in the ambient metric. Choose a relatively
compact ambient neighborhood of $\Sigma$ and a smaller compact safety
neighborhood. Fix $r_*>0$ much smaller than the distance separating
their boundaries and the injectivity radii on the larger compact
neighborhood. Short paths used in the proof then stay in the region
of controlled smooth geometry, and the squared-distance functions
used for pairs less than $4r_*$ apart are smooth. No completeness of
an auxiliary ambient extension is needed.

\paragraph{Regular distance levels.}
If $0<\rho(x)<r_*$, every nearest point of $\Sigma$ to $x$ belongs to
$\Lambda$. Indeed, a shortest segment from $x$ to the full boundary
stays in $\Omega$ until its endpoint, and a shorter terminal subsegment
supplies an interior tangent ball there. The nearest-point set is
compact. It therefore lies in the interior of a compact smooth
hypersurface patch $Q\subset\Lambda$ with smooth boundary. Choose
$Q$ sufficiently close to that set that
$\dist(y,q)^2$ is smooth for $q\in Q$ and $y$ near $x$.

Because $\Lambda$ is embedded, there is an ambient open set
$\mathcal V$ containing the nearest-point set with
$\mathcal V\cap\Lambda\subset\operatorname{int}_{\Lambda}Q$.
After shrinking the neighborhood of $x$, every nearest point to each
such $y$ lies in $\mathcal V$. Otherwise, compactness of
$\Sigma\setminus\mathcal V$ would give a nearest point of $x$
outside $\mathcal V$. These nearby nearest points belong to
$\Lambda$ by the preceding tangent-ball argument, and hence lie
in the interior of $Q$. Thus
\[
 \rho(y)^2=\min_{q\in Q}\dist(y,q)^2
\]
on an open neighborhood of $x$. Apply
\cite[Theorem~3 and the compact-parameter-with-boundary extension
in the proof of Theorem~1]{Rifford2004}. A countable family of these
neighborhoods covers the small distance collar, so the critical
values of $\rho^2$ form a null set. On each compact interval of
positive distances, the square-root map is Lipschitz and
$\partial_C(\rho^2)=2\rho\,\partial_C\rho$, where $\partial_C$
denotes the Clarke differential. Thus the positive critical values
of $\rho$ also form a null set. Choose $s_i\downarrow0$ such that
\[
 0\notin\partial_C\rho(x)\qquad\text{whenever }\rho(x)=s_i.
\]
The compact parameter patch is the reason this theorem applies even
though the accessible locus $\Lambda$ need not be closed.

\paragraph{Positive reach of the superlevels.}
The distance function is locally semiconcave away from $\Sigma$.
To obtain smooth upper supports, take a minimizing segment from
$\Sigma$ to a point of the collar and retain a short terminal piece.
Distance from its earlier endpoint, plus the length of the preceding
piece, is a smooth upper support near the terminal point. The local
Hessian comparison estimate gives a uniform upper Hessian bound on
a sufficiently small fixed ball. Consequently, in a geodesic
coordinate ball,
\begin{equation}\label{eq:rpi:smoothing-semiconcavity}
 \rho(y)\le\rho(x)+\langle p,\log_x y\rangle
                  +C\dist(x,y)^2
 \qquad(p\in\partial_C\rho(x)).
\end{equation}
Here the superdifferential of the semiconcave function agrees with its
Clarke differential.

We include the regular-superlevel implication underlying the
positive-reach criterion associated with \cite{Bangert1982}.
At a point of a regular level, let
$\mathcal K=\partial_C\rho(x)$. This is a compact convex set not
containing zero. Separation gives a direction $v_0$ such that
$\langle p,v_0\rangle>0$ for every $p\in\mathcal K$. The directional
derivative of a semiconcave function is
$\min_{p\in\mathcal K}\langle p,v\rangle$. Every direction for which
this minimum is positive therefore enters the superlevel. Separation
of convex cones implies that the polar of its contingent tangent
cone, and in particular each proximal outward normal, is contained
in $-\operatorname{cone}(\mathcal K)$. Since $\mathcal K$ is convex,
each unit proximal outward normal can be written
$\mathbf n=-p/|p|$ with $p\in\mathcal K$.

Compactness of the level and upper semicontinuity of the
superdifferential give a uniform lower bound $|p|\ge\eta>0$ there.
For a nearby point $y$ of the superlevel,
Equation~\eqref{eq:rpi:smoothing-semiconcavity} yields
\[
 \langle\mathbf n,\log_x y\rangle
       \le\frac C\eta\dist(x,y)^2.
\]
This gives uniqueness of sufficiently nearby nearest-point
projection. Indeed, suppose that $z$ had two distinct nearest
points $x,y$, at the same distance $a$. In a sufficiently small
bounded-geometry ball, the Taylor expansion of squared distance
from $z$ gives
\[
 0=\dist(z,y)^2-\dist(z,x)^2
 \ge-2a\langle\mathbf n,\log_x y\rangle
                         +\frac12\dist(x,y)^2.
\]
For $a<\eta/(4C)$ this contradicts the preceding bound. Enlarging
the constants slightly covers the curvature terms. Compactness
excludes distant distinct nearest points when $a$ is small and
makes this projection radius uniform. Thus
\[
 F_i=\{x\in\Omega:\rho(x)\ge s_i\}
\]
has positive reach. It is closed in the ambient manifold because
$\Sigma$ is the full boundary and $s_i>0$, and its boundary is
compact. The same separating direction $v_0$ shows that its outward
normal cone is pointed: it cannot contain opposite nonzero normals.
Only the regular-superlevel implication has been used.

\paragraph{Offsets and their curvature.}
Fix a regular value $s=s_i$ and put $F=F_i$. For
$0<\tau<\operatorname{reach}(F)$, sufficiently small compared with
$s$ and the ambient injectivity scale, define
\[
 O_\tau=\{x:\dist(x,F)<\tau\},\qquad
 \Gamma_\tau=\partial O_\tau.
\]
The closure of $O_\tau$ lies in $\Omega$, since
$\rho\ge s-\tau>0$ there. Moreover,
$\Omega\setminus O_\tau$ lies in the $s$-collar of $\Sigma$.

Each point of $\Gamma_\tau$ has a unique footpoint $q\in F$ and
can be written $p=\exp_q(\tau v)$, where $v$ is a unit outward
normal to $F$. The same ray realizes distance to $F$ for every
positive parameter below $\min\{\operatorname{reach}(F),r_*\}$.
These facts follow from uniqueness of
projection, first variation and continuation of a minimizing normal
segment; they are also the nearest-point and normal-cone facts used
in \cite[Section~3.1]{BiZhuSmoothing2026}. No differentiability
of the projection is required.

The normal ray supplies tangent balls on both sides of
$\Gamma_\tau$. The interior supporting radius is bounded below
when $\tau$ stays away from zero. For the exterior supporting ball,
fix $0<R<\min\{\operatorname{reach}(F),r_*\}$ first and take
$\tau<R/2$. Continuing the normal ray to $R$ gives an exterior
supporting radius $R-\tau\ge R/2$. Local graph comparison with
these spheres gives a compact $C^{1,1}$ hypersurface.
For this fixed $F$, the principal curvatures have a uniform lower
bound as $\tau\downarrow0$; their upper bound may grow like
$\tau^{-1}$.

The offset is now a $C^{1,1}$ hypersurface lying strictly inside
$\Omega$. To finish its construction we must transfer the positive
curvature of accessible boundary points to this offset. We do so
by comparison with smooth interior test domains.
Let a smooth region $D\subset O_\tau$ touch $\Gamma_\tau$ from
inside at $p=\exp_q(\tau v)$. Its outward normal at contact is the
parallel translate of $v$ along the normal ray.

Suppose first that the unit normal cone of $F$ at $q$ is a singleton.
Take a shortest segment from $q$ to $\Sigma$. Its initial direction
is an outward normal to $F$: its negative is the gradient of a
smooth distance upper support, and the superlevel inequality places
the initial direction in the polar tangent cone. It must therefore
equal $v$. Its endpoint $z\in\Sigma$ is accessible and smooth.
Figure~\ref{fig:rpi:smoothing-offset} shows the order of the two
offsets along this segment.
Enlarging $D$ by distance $s-\tau$ touches $\Sigma$ at $z$. The
enlargement stays on the $\Omega$ side, since every point of $D$
has distance at least $s-\tau$ from $\Sigma$.

\begin{figure}[tbp]
\centering
\begin{tikzpicture}[x=1.12cm,y=1cm,>=Latex,font=\small]
 \fill[blue!4] (0,-.55) rectangle (7,-.05);
 \fill[blue!11] (2.7,-.55) rectangle (7,-.05);
 \fill[blue!21] (4.7,-.55) rectangle (7,-.05);
 \draw[gray!65,dashed] (0,-.7)--(0,.7);
 \draw[blue!60!black,dashed] (2.7,-.7)--(2.7,.7);
 \draw[blue!60!black,dashed] (4.7,-.7)--(4.7,.7);
 \draw[thick] (0,0)--(4.7,0);
 \fill (0,0) circle (1.8pt);
 \fill (2.7,0) circle (1.8pt);
 \fill (4.7,0) circle (1.8pt);
 \node[above left] at (0,0) {$z$};
 \node[above right] at (2.7,0) {$p$};
 \node[above right] at (4.7,0) {$q$};
 \node[below] at (0,-.74) {$\Sigma$};
 \node[below] at (2.7,-.74) {$\Gamma_\tau=\partial O_\tau$};
 \node[below] at (4.7,-.74) {$\partial F$};
 \node at (1.25,-.3) {$\Omega$};
 \node at (3.7,-.3) {$O_\tau$};
 \node at (5.85,-.3) {$F$};
 \draw[<->] (0,1.1)--node[above] {$s$}(4.7,1.1);
 \draw[<->] (2.7,.68)--node[above] {$\tau$}(4.7,.68);
 \draw[->,thick] (2.7,.25)--(1.85,.25)
       node[midway,above] {$\mathbf n_{O_\tau}$};
 \draw[->,thick] (4.7,.25)--(3.95,.25)
       node[midway,above] {$v$};
\end{tikzpicture}
\caption{The offset construction along a minimizing segment in the
singleton-normal case. First take $F=\{\rho\ge s\}$, then enlarge it
by $\tau<s$ to $O_\tau$; the gap to $\Sigma$ on this segment is
$s-\tau$. The outward normals point toward $\Sigma$, hence equal
$-\nu$ in the exterior application. This local schematic asserts no
tubular coordinates at singular points.}
\label{fig:rpi:smoothing-offset}
\end{figure}

Let $C_R\ge0$ bound the negative part of the ambient Ricci curvature
on the fixed neighborhood. Along the touching geodesic, the normal
Riccati equation gives
\[
 H(\text{enlarged }D\text{ at }z)
        \le H(D\text{ at }p)+C_R(s-\tau).
\]
First-contact comparison with the accessible part of $\Sigma$ gives
$H(\text{enlarged }D\text{ at }z)\ge\gamma$.

Here is the endpoint argument if the enlargement is focal at $z$.
Localize $D$ near $p$ and shrink it strictly away from $p$, preserving
the touching segment, so that $p$ is the unique terminal footpoint.
In tangent coordinates with outward normal in the positive vertical
direction, replace its boundary graph $u$ by
$u-a|x|^2$, where $a>0$ can be arbitrarily small. This shrinks the
domain and increases its outward second fundamental form at contact
by $2a$ times the tangent metric. The boundary index form of the
original minimizing normal segment is nonnegative; the perturbation
adds $2a|J(0)|^2$ for each variation field $J$. Variations with
$J(0)=0$ have positive index because the segment is shorter than
the ambient conjugacy scale. Thus the perturbed segment is nonfocal
through its endpoint. Apply the Riccati comparison to the perturbed
domain and let $a\downarrow0$. Its mean curvature at contact tends
to that of $D$, proving also in this case that
\[
 H(D\text{ at }p)\ge\gamma-C_Rs.
\]

Suppose next that the normal cone contains a direction not parallel
to $v$. Convexity supplies a nontrivial one-sided angular arc through
$v$. Exponentiating this arc at radius $\tau$ produces a curve in
$\Gamma_\tau$. The test boundary has the same first derivative at
contact. Second-order comparison along the arc, even when it is
one-sided, therefore gives a principal curvature at least
$1/(2\tau)$ for small $\tau$. The exterior supporting ball gives
a lower bound $-C_F$ for all the other principal curvatures,
independent of $\tau$ for this fixed $F$. Hence
\[
 H(D\text{ at }p)\ge\frac1{2\tau}-(n-2)C_F.
\]
This exceeds any fixed lower bound when $\tau$ is sufficiently
small. The pointedness proved above excludes the remaining apparent
possibility of opposite normals without an angular arc.

Choose the regular value $s<\varepsilon$ sufficiently small that
$C_Rs<\varepsilon/4$. Then choose $\tau$ sufficiently small compared
with $s$, the reach of this fixed $F$, and its curvature constants.
The resulting $C^{1,1}$ hypersurface has outward mean curvature at
least $\gamma-\varepsilon/4$ in the barrier sense. At almost every
twice-differentiability point of a $C^{1,1}$ graph, quadratic test
functions give the same inequality for its classical almost-everywhere
mean curvature. The order of these choices is important: neither
the reach nor $C_F$ is asserted to remain uniform as $s\downarrow0$.

\paragraph{Smooth graph approximation.}
We use the following direct fact. A compact cooriented embedded
$C^{1,1}$ hypersurface with almost-everywhere outward mean curvature
at least $c$ has smooth embedded approximations converging in $C^1$
whose mean curvature is at least $c-o(1)$.

To prove it, choose a smooth ambient vector field $V$ uniformly
transverse to the hypersurface, by smoothly approximating its
continuous unit normal near the compact hypersurface. The $C^1$
inverse function theorem and compactness give a fixed flow tube for
$V$, in which each sufficiently short flow line meets the hypersurface
exactly once. Local graph convolution and a partition of unity give
a smooth $C^1$ approximation $\Gamma_0$ sufficiently close to lie in
this tube and remain uniformly transverse to $V$. Projection along
these fixed flow lines is a global diffeomorphism. The original
hypersurface is therefore the flow graph of a $C^{1,1}$ function
$u$ over $\Gamma_0$. This fixed transverse flow avoids any assumption
about the normal injectivity radius of an arbitrary $C^1$ approximation.

In local coordinates its mean-curvature operator has the form
\[
 H(u)=A^{ab}(x,u,Du)\,\partial_{ab}u+B(x,u,Du),
\]
where the coefficients are smooth on a compact set of allowed graph
jets. The sign of $A$ depends on the graph convention and is immaterial
for the following averaging argument. Choose a finite smooth atlas
with subordinate partition of unity $\chi_\alpha$ and smooth $u$
using normalized nonnegative convolution kernels
$\rho_{\alpha,\zeta}$. The functions $u_\zeta$ converge to $u$ in
$C^1$ and have uniformly bounded second derivatives. Their second
derivatives equal the partition-weighted averages of $D^2u$ plus
a uniformly vanishing term. Indeed, differentiating the partition
introduces differences $u_\zeta-u$ and $Du_\zeta-Du$, which tend
uniformly to zero, while the sums of the derivatives of the partition
vanish.

Because $D^2u$ is uniformly bounded, the coefficient replacement
inside those averages,
\[
 A(x,u_\zeta,Du_\zeta)\longmapsto A(y,u(y),Du(y)),
\]
gives a uniformly vanishing commutator. The
lower-order coefficients converge uniformly as well. Coordinate
changes and connection terms are included in $B$; their additional
smoothing errors are controlled by the same $C^1$ convergence.
Consequently
\[
 H(u_\zeta)(x)
 =\sum_\alpha\chi_\alpha(x)
    \int\rho_{\alpha,\zeta}(x,y)H(u)(y)\,dy+o(1)
 \ge c-o(1)
\]
uniformly in $x$, proving the asserted approximation.

Apply this construction to the fixed hypersurface $\Gamma_\tau$.
Its positive distance from $\Sigma$ ensures that a sufficiently
small $C^1$ perturbation remains inside $\Omega$. Choose the region
on the same side of the smoothed boundary that contains the deep
part of $\Omega$. It differs from $O_\tau$ only in a tubular
neighborhood of $\Gamma_\tau$. Choose the width of this neighborhood
smaller than both its distance from $\Sigma$ and
$\varepsilon-s$. The removed collar is then contained in the
$\varepsilon$-collar of $\Sigma$. Finally choose $\zeta$ so small
that the mean-curvature loss is less than $\varepsilon/2$. The
resulting region is $\Omega_\varepsilon$, and its mean curvature is
at least $\gamma-\varepsilon$. This proves the proposition.
\end{proof}

For the almost-minimizing frontiers used here, the accessible locus
is exactly the regular part. At a point touched by a ball contained
in the exterior, a tangent cone lies on one side of a plane. The
half-space property for minimizing boundaries and small-excess
regularity make the point regular. Conversely, a regular point has
a sufficiently small supporting ball on the chosen side. Thus the
smooth embeddedness and local one-sided hypotheses of
Proposition~\ref{prop:rpi:one-sided-smoothing} hold. The proposition
is applied to the full compact frontier. An application to only
selected boundary components would also require a positive distance
from the remaining components; compactness of their union alone
does not imply that separation.

\paragraph{Application to a detached frontier.}
Let $E'$ be a detached minimizing exterior from the signed
variational construction, with $H_\nu\le-\gamma$ on its regular
frontier. Apply the proposition to $\Omega=E'$.
Its outward normal at the inner frontier is $-\nu$, so
$H_{\mathrm{out},\Omega}=-H_\nu\ge\gamma$. With
$\varepsilon=j^{-1}$, write $\Omega_j=\Omega_{j^{-1}}$. Since
$\Omega_j\subset\Omega$, the new frontier lies toward the
asymptotically flat end. With the end-oriented normal it satisfies
\begin{equation}\label{eq:rpi:smoothing-sign}
 H_\nu(\partial\Omega_j)\le-\gamma+j^{-1}<0
\end{equation}
for large $j$. Its compact removed region contains the original
removed region. In the variational definition of capacity this makes
the admissible class smaller, and hence
\begin{equation}\label{eq:rpi:capacity-monotonicity}
 \mathfrak c(\partial\Omega_j,h)\ge\mathfrak c(\Sigma,h).
\end{equation}
If $\Omega_j$ has bounded components, retain only the component
containing the asymptotic end. This selects whole components of its
smooth compact boundary, preserves the curvature sign, and enlarges
the conductor further. The same capacity comparison therefore holds
for the resulting connected smooth exterior. The argument uses only
the shrinking collar, the curvature sign and conductor inclusion;
no boundary-area convergence is needed.

This smooth exterior is complete with its boundary included. For
each fixed $\varepsilon$, the harmonic perturbation is uniformly
comparable with the starting metric, and the retained closed
exterior lies a positive distance from the original frontier.
The finite-length concatenation argument in
Lemma~\ref{lem:g-detachment} therefore gives completeness for its
intrinsic length distance. These are the smooth boundary and
completeness properties needed in the mass--capacity lemma below.

\section{Positive mass and the mass--capacity inequality}
\label{sec:rpi:smooth-repair}

The preceding construction reduces the singular-frontier argument
to a smooth mass--capacity inequality. We prove that inequality
from a positive mass theorem, independently of
Theorem~\ref{thm:rpi}. The conformal-flow argument at this stage
has a harmonically flat end, and that is the class we need. A harmonically flat
end has metric $U^{4/(n-2)}\delta$, with $U>0$ Euclidean harmonic and
$U\to1$. We use the independently stated smooth positive mass result
\cite[Corollary~1.6]{BrendleWang2026}. Its actual hypotheses are strict
positive scalar curvature and an all-order end expansion
\begin{equation}\label{eq:rpi:bw-end}
 g=(1+\alpha r^{2-n})\delta+O_j(r^{2-n-\sigma}),\qquad \sigma>0,
\end{equation}
for every fixed $j$. The complete smooth manifold in that corollary
and its governing definitions has no orientability assumption. In
dimensions at least four other ends are permitted. Its mass is
$(n-1)\alpha$, whereas our ADM mass is
$(n-2)\alpha/2$; their signs agree. The following argument uses it only
for $n\ge8$; the lower-dimensional numerical positive mass theorem
and \cite{BrayLee2009} provide the classical alternative.

The dimension-descent proof of the cited theorem uses a singular-set
covering estimate for weighted perimeter minimizers with a volume
forcing. Appendix~\ref{sec:rpi:weighted-minimizers} proves the
needed estimate for that class, using
\cite[Theorem~1.3]{NaberValtorta2020}; its role in
\cite[Definition~3.19, Lemma~3.20, and Theorem~3.34]{BrendleWang2026}
is specified there. With that input accounted for, the remaining
steps of the mass--capacity proof concern smooth metrics: first
pass from strict to nonnegative scalar curvature on harmonic ends,
then double the smooth exterior and repair its seam.

\subsection{Nonnegative mass on harmonic ends}

\begin{lemma}\label{lem:rpi:harmonic-pmt}
A smooth complete boundaryless manifold of dimension $n\ge8$ with finitely many
ends, all harmonically flat, and nonnegative scalar curvature has nonnegative
ADM mass at each end.
\end{lemma}
\begin{proof}
Work on the connected component containing the chosen end.
Choose a positive smooth function $\rho$, equal to $r^{-n-\beta}$ on
each distant end, with $0<\beta<1$. The global Sobolev inequality on a
manifold with finitely many Euclidean ends, followed by exhaustion,
gives the positive solution of
\[
 -\Delta v=\rho,\qquad v\longrightarrow0\quad\hbox{on every end}.
\]
More explicitly, $\rho\in L^{2n/(n+2)}$, so the energy completion of
compactly supported functions gives a weak solution by the coercive
Dirichlet form. Positivity follows by testing the negative part.
The local elliptic estimates and end barriers give a bounded smooth
solution; the homogeneous decaying maximum principle gives uniqueness.
On each harmonically flat end, $Uv$ satisfies a Euclidean Poisson
equation with right side $-U^{(n+2)/(n-2)}\rho$. The Newton potential,
split into inner, comparable, and outer annuli, gives
\[
 v=A r^{2-n}+O_j(r^{2-n-\beta'})
 \quad\text{for }0<\beta'<\beta.
\]
Equivalently, subtract the Newton kernel's monopole and use the
finite $\beta'$ moment of the right-hand side. The contribution of
the finite inner boundary is a decaying harmonic function with the
same type of expansion. Scaled interior estimates give all fixed
orders of derivatives. Thus $(1+t v)^{4/(n-2)}g$ satisfies
Equation~\eqref{eq:rpi:bw-end}, is complete, and has strictly positive
scalar curvature for every $t>0$, since
\[
 -\frac{4(n-1)}{n-2}\Delta(1+t v)+R_g(1+t v)>0.
\]
Apply the cited positive mass theorem and let $t\downarrow0$.
The ADM mass changes by $2tA$ at the chosen end, which proves the
claim.
\end{proof}

\subsection{Doubling and conformal repair}

We can now turn the capacity of a smooth boundary into a mass
decrease on a complete manifold. The doubling must keep its second
end complete while positive mass is applied; a positive parameter
will do this, and it can be sent to zero afterward.

The following doubling, seam smoothing, and conformal repair have
their three-dimensional antecedent in
\cite[Definition~17 and Theorem~9, pp.~206--211]{Bray2001}.
The normalization there gives the same three-dimensional capacity
bound. The higher-dimensional positive-mass input is instead the
smooth theorem and the reduction just established.

\begin{lemma}[Smooth mass--capacity inequality]
\label{lem:rpi:mass-capacity}
Let $(N,h)$, of dimension $n\ge8$, be a smooth connected exterior,
complete including its compact smooth boundary, with exactly one end
and compact end complement. Suppose the end is harmonically flat
and the scalar curvature is nonnegative. Write $\Gamma$ for its
boundary. If $H_\nu\le0$ for the normal toward the end, then
$m_{\rm ADM}(h)\ge\mathfrak c(\Gamma,h)$.
\end{lemma}
\begin{proof}
If $\Gamma$ is empty, its capacity is zero and the assertion is
Lemma~\ref{lem:rpi:harmonic-pmt}. Assume henceforth that $\Gamma$
is nonempty. Let $\phi=1-w$, where $w$ is the potential in
Equation~\eqref{eq:rpi:capacity}, and write $c=\mathfrak c(\Gamma,h)$.
On the doubled manifold define $\psi=\phi$ on the first copy and
$\psi=-\phi$ on the second. In signed normal distance across the seam,
$\psi$ and its first derivative agree: its value is zero and both
normal derivatives equal $\partial_\nu\phi$. For $0<a<1$, put
\[
 V_a=\frac{1+a+(1-a)\psi}{2},\qquad
 h_a=V_a^{4/(n-2)}\bar h,
\]
where $\bar h$ is the doubled corner metric. The factor is bounded
below by $a$, tends to one at the first end and to $a$ at the second,
and is weakly harmonic. Consequently $h_a$ is complete, has two
harmonically flat ends after constant rescaling at the second, and has
nonnegative scalar curvature away from the seam.

Choose the common transverse normal from the second copy to the first.
The two mean curvatures of $\bar h$ with this normal are $-H_\nu$ and
$H_\nu$. Their jump in the corner condition is therefore $-2H_\nu\ge0$.
The normal derivative of $V_a$ is continuous, so its contribution
cancels in the conformal jump formula; the jump for $h_a$ has the same
sign. We now explain the analytic smoothing step to avoid extending
the dimensional or topological statement of a corner theorem without
justification.

Apply the local Gaussian-collar construction of
\cite[Section~3, Proposition~3.1 and Equation~(36)]{Miao2002}.
Its tensor convolution uses the coorientation of this collar and
is independent of the topology or number of ends outside it.
It gives uniformly comparable $C^2$ metrics, equal to $h_a$ off
a shrinking collar, whose negative scalar curvature has a uniform
bound. For each fixed collar parameter, approximate this metric
in $C^2$ by a smooth metric, keeping it unchanged outside a slightly
larger collar. Scalar curvature is continuous in the metric's
$C^2$ jets, so choosing the approximation sufficiently close
preserves the negative-scalar bound up to a fixed enlargement.
Denote these smooth metrics by $h_{a,\delta}$.
Their negative scalar curvature is supported in a collar whose
volume tends to zero. Choose a smooth
$q_\delta\ge(R_{h_{a,\delta}})_-$, uniformly bounded and supported in
a slightly larger shrinking collar. Solve
\[
 \Delta_{h_{a,\delta}} z_\delta+
       \frac{n-2}{4(n-1)}q_\delta z_\delta=0,
 \qquad z_\delta\longrightarrow1\quad\text{at both ends}.
\]
Here is the required existence and mass estimate on this two-ended
manifold. Put $a_n=4(n-1)/(n-2)$, $2^*=2n/(n-2)$ and
$p=2n/(n+2)$. A fixed complete metric with these ends satisfies
$\|f\|_{2^*}^2\le S_a\int|df|^2$: combine the Euclidean end
inequalities with a compact-core Poincar\'e estimate having trace on
one fixed end annulus. Uniform metric comparison gives the same
inequality for $h_{a,\delta}$, uniformly in small $\delta$.
Since $\|q_\delta\|_{n/2}\to0$, the form
\[
 B_\delta(f,f)=a_n\int|df|^2-\int q_\delta f^2
       \ge\frac{a_n}{2}\int|df|^2
\]
is coercive on the homogeneous energy completion of compactly
supported smooth functions. Solving
$B_\delta(u_\delta,f)=\int q_\delta f$ gives
\[
 \|u_\delta\|_{2^*}+\|du_\delta\|_2
       \le C_a\|q_\delta\|_p.
\]
Testing the negative part gives $u_\delta\ge0$. Local elliptic
regularity and the harmonic exterior equation give a smooth solution
$z_\delta=1+u_\delta\ge1$, tending to one on both ends; the end
limits also follow from $u_\delta\in L^{2^*}$.
The corrected smooth metric
$z_\delta^{4/(n-2)}h_{a,\delta}$ has nonnegative scalar curvature.
It is still harmonically flat at both ends. In separately normalized
end coordinates write
$u_\delta=A_{i,\delta}r_i^{2-n}+O_j(r_i^{1-n})$;
the coefficients are nonnegative because $u_\delta\ge0$.
Flux integration over both ends gives
\[
 a_n(n-2)\omega\sum_{i=1}^2A_{i,\delta}
   =\int q_\delta z_\delta\,dV
   \le\|q_\delta\|_1+C_a\|q_\delta\|_p^2=o(1).
\]
Thus each coefficient tends to zero, and each end mass converges to
that of $h_a$.
Lemma~\ref{lem:rpi:harmonic-pmt} proves nonnegativity of this limiting
mass. This argument is local at the seam and accommodates both ends;
it does not invoke the one-ended global statement of
\cite[Theorem~1]{Miao2002} outside its stated class.

On the retained end,
\[
 \phi=1-cr^{2-n}+O_j(r^{1-n}),\qquad
 V_a=1-\tfrac12(1-a)cr^{2-n}+O_j(r^{1-n}),
\]
so the conformal ADM formula gives
\[
 0\le m_{\rm ADM}(h_a)=m_{\rm ADM}(h)-(1-a)c.
\]
Let $a\downarrow0$. Both ends remained complete for every application
of the smooth positive mass theorem; no theorem about a compactified
point was used.
\end{proof}

The signed variational construction, one-sided smoothing and
Lemma~\ref{lem:rpi:mass-capacity} now give the smooth comparison
needed for a singular time slice: an enclosing smooth conductor has
at least the original capacity, while its end mass is unchanged by
moving the frontier. Section~\ref{sec:rpi:flow-mass-capacity} applies
this comparison to a small harmonic perturbation and removes that
perturbation. The same positive two-ended construction will provide
the complete inputs used in Section~\ref{sec:rpi:lee-completion}.
